\documentclass[12pt]{article}
% Préambule commun à tous les documents extraits de « La Revue CAPES »
\usepackage[T1]{fontenc}
\usepackage[utf8]{inputenc}
\usepackage{lmodern}
\usepackage{amsmath,amssymb,amsthm,amsfonts,mathrsfs}
\usepackage{setspace}
\usepackage{listings}
\usepackage{pgfplots}
\pgfplotsset{compat=1.18}
\usepackage{longtable}
\usepackage{adjustbox}
\usepackage{lettrine}
\usepackage{tkz-tab}
\usepackage{changepage}
\usepackage{pdflscape}
\usepackage{graphicx}
\usepackage{tikz}
\usepackage{xcolor}
\usepackage[a4paper,top=2cm,bottom=2.5cm,left=2.5cm,right=2cm]{geometry}
\usepackage{fancyhdr}
\usepackage{draftwatermark}
\SetWatermarkText{La RevueCapes}
\SetWatermarkLightness{0.9}
\SetWatermarkScale{2}
\usepackage{hyperref}
\hypersetup{colorlinks=true,linkcolor=blue!50!black,urlcolor=blue!50!black}

\definecolor{revue}{RGB}{20,60,120}

\pagestyle{fancy}
\fancyhf{}
\renewcommand{\headrulewidth}{0.4pt}
\setlength{\headheight}{15pt}

% \entete{Matière}{Titre}{Type de document}
\newcommand{\entete}[3]{%
  \fancyhead[L]{\small\textsc{La Revue CAPES} -- #1}%
  \fancyhead[R]{\small #2}%
  \fancyfoot[C]{\thepage}%
  \thispagestyle{plain}%
  \begin{center}
    {\color{revue}\rule{\textwidth}{1.2pt}}\\[4pt]
    {\small\textsc{Concours directs CAP-CEG / CAPES -- Burkina Faso}}\\[6pt]
    {\LARGE\bfseries\color{revue} #2}\\[6pt]
    {\large\itshape #1 -- #3}\\[2pt]
    {\color{revue}\rule{\textwidth}{1.2pt}}
  \end{center}
  \vspace{0.5cm}%
}

% Encadré pour les remarques ajoutées dans les corrigés
\newcommand{\remarque}[1]{\par\medskip\noindent\fcolorbox{revue}{revue!6}{\parbox{\dimexpr\linewidth-2\fboxsep-2\fboxrule}{\textbf{\color{revue}Remarque.} #1}}\par\medskip}


\begin{document}
\entete{Culture générale}{Corrigé -- Session 2017}{Correction}
\begin{spacing}{1.5}

\textbf{Introduction}

Dans une société où la réussite est souvent perçue comme le seul critère de valorisation, il est facile d'oublier que l'effort en lui-même constitue une richesse. Gandhi, dans cette pensée, invite à repenser notre relation à l'effort et au succès : il affirme que la satisfaction réside moins dans l'atteinte des objectifs que dans la sincérité et l'intensité de l'engagement. Il nous rappelle également que l'effort constant en quête de perfection est une obligation morale qui porte en elle sa propre récompense. Cette réflexion soulève des interrogations fondamentales : pourquoi l'effort aurait-il plus de valeur que le résultat ? En quoi cette quête d'amélioration continue constitue-t-elle une victoire personnelle et universelle ? Après avoir clarifié cette pensée, nous discuterons de sa pertinence et de ses limites dans le contexte contemporain.


\textbf{Développement}

\textbf{I. La valeur intrinsèque de l'effort}

Gandhi souligne que l'effort, en tant qu'acte volontaire et sincère, porte en lui-même une satisfaction profonde. Cette idée repose sur le fait que l'effort est une manifestation de notre volonté de progresser et de nous améliorer, indépendamment des résultats obtenus.  

\textbf{1. Une source de satisfaction personnelle} 
 
   L'effort déployé dans un travail ou une activité procure une forme d'accomplissement. Il renforce la confiance en soi et donne un sens à nos actions, même lorsque le succès n'est pas garanti. Par exemple, un étudiant qui se consacre avec sérieux à ses études tire une fierté de son engagement, quelle que soit l'issue de ses examens.  
   
\textbf{2. Une éthique universelle 
   }
   
   Dans de nombreuses cultures et traditions, l'effort est perçu comme une vertu morale. Les philosophies stoïcienne ou bouddhiste, par exemple, prônent l'idée que l'essentiel est dans la maîtrise de soi et dans le travail sur ses imperfections, plutôt que dans la recherche de récompenses extérieures.

\textbf{II. L'effort comme une victoire en soi}

Gandhi affirme que le plein effort équivaut à une victoire. Cette vision valorise la démarche plus que le résultat, et encourage un dépassement de soi constant.  
\textbf{1. Le dépassement personnel
}
   Faire de son mieux, quel que soit le domaine, permet de repousser ses limites et de progresser. Ce processus renforce le caractère et l'endurance. Par exemple, un athlète qui s'entraîne sans relâche pour une compétition peut considérer son engagement comme une victoire, même s'il ne gagne pas la médaille d'or. 
    
\textbf{2. Une victoire intérieure}
   La satisfaction tirée de l'effort sincère va au-delà des jugements externes. Elle permet de vivre en accord avec ses valeurs, ce qui constitue une forme de victoire personnelle. Gandhi, dans sa propre vie, a illustré cette idée en luttant pour l'indépendance de l'Inde avec des moyens non violents, malgré les immenses défis auxquels il était confronté.

\textbf{III. Les limites et les enjeux contemporains de cette pensée}

Si la philosophie de Gandhi est inspirante, elle peut sembler difficilement applicable dans certains contextes modernes.
 
\textbf{1. L'importance des résultats dans la société actuelle}
  
   Dans un monde orienté vers la performance et les résultats mesurables, il est parfois difficile de valoriser uniquement l'effort. Par exemple, dans un cadre professionnel, un travailleur pourrait déployer beaucoup d'efforts, mais ne pas être reconnu s'il n'atteint pas les objectifs fixés. 
    
\textbf{2. Le risque de l'épuisement  }

   Exiger un effort constant sans se soucier des résultats peut aussi conduire à une pression excessive et à un épuisement moral ou physique. Il est donc nécessaire de trouver un équilibre entre persévérance et repos, entre quête de perfection et acceptation de ses limites.  

\textbf{IV. L'effort comme chemin vers une perfection morale}

Enfin, Gandhi élève l'effort au rang de devoir universel. Selon lui, chercher continuellement à s'améliorer est un moyen de se réaliser pleinement et de contribuer au bien commun.
  
\textbf{1. Un devoir moral } 
   L'effort constant incarne une quête de perfection morale et d'éthique. En travaillant sur soi-même, on devient plus apte à servir les autres et à avoir un impact positif sur la société.  
   
\textbf{2. La récompense intérieure  }

   Cet effort ne dépend pas de l'approbation des autres ni de la reconnaissance extérieure, mais de la satisfaction personnelle et du sentiment d'avoir fait ce qui est juste. Cela est particulièrement pertinent dans un monde où la recherche de validation extérieure peut devenir oppressante.


\textbf{Transition vers la conclusion}

Ainsi, la pensée de Gandhi nous invite à repenser notre rapport à l'effort, à voir dans le cheminement un enrichissement personnel et une victoire en soi. Pourtant, cette vision idéale nécessite d'être nuancée pour rester applicable aux réalités contemporaines, où effort et résultats doivent souvent être conciliés. 



\quad

\quad

\textbf{Conclusion}

La pensée de Gandhi nous enseigne que l'essence de la satisfaction et de la victoire réside dans l'effort sincère et constant plutôt que dans la simple quête de résultats. Elle valorise une vision humaniste où l'amélioration personnelle et le dépassement de soi priment sur la réussite matérielle ou sociale. Cette approche, bien qu'exigeante, nous invite à trouver un sens profond dans nos actions quotidiennes et à embrasser l'idée que la perfection n'est pas un objectif fixe, mais un cheminement continu. Cependant, dans un monde où les résultats concrets sont souvent indispensables, il est nécessaire de trouver un équilibre entre cette philosophie idéale et les réalités pragmatiques de la vie moderne. Gandhi nous propose ainsi une leçon intemporelle : c'est dans la sincérité de nos efforts que réside notre véritable accomplissement.

\quad
\quad

\end{spacing}
\end{document}
